\begin{trapbox}
	\begin{description}[style=nextline, leftmargin=2.8em, labelsep=0.5em, itemsep=0.35em]
		\item[\textbf{\textcolor{CTrap}{易错点一（混淆条件与交）：}}] 将$P(B \mid A)$等同于$P(AB)$，认为条件概率就是“$A$且$B$”的概率。
		\item[\textbf{\textcolor{CTrap}{正确理解（区分公式意义）：}}] $P(B \mid A)$表示在已知$A$发生的条件下$B$的概率，$P(AB)$表示$A$与$B$同时发生的概率，两者一般不同。
		\item[\textbf{\textcolor{CTrap}{错因分析（忽略条件背景）：}}] 学生忽略了“已知$A$发生”这一前提，直接把条件事件当作交事件。
	\end{description}

	\medskip
	\begin{description}[style=nextline, leftmargin=2.8em, labelsep=0.5em, itemsep=0.35em]
		\item[\textbf{\textcolor{CTrap}{易错点二（忽视分母为零）：}}] 在$P(A)=0$时仍然使用公式$P(B \mid A)=\frac{P(AB)}{P(A)}$，导致无意义或错误结论。
		\item[\textbf{\textcolor{CTrap}{正确理解（验证非零条件）：}}] 条件概率定义的前提是$P(A)>0$，否则条件概率无定义。
		\item[\textbf{\textcolor{CTrap}{错因分析（忘记检查条件）：}}] 学生使用公式前未检查条件事件概率是否为正，直接套用。
	\end{description}

	\medskip
	\begin{description}[style=nextline, leftmargin=2.8em, labelsep=0.5em, itemsep=0.35em]
		\item[\textbf{\textcolor{CTrap}{易错点三（颠倒条件先后）：}}] 误认为$P(B \mid A)=P(A \mid B)$，从而交换条件与结论。
		\item[\textbf{\textcolor{CTrap}{正确理解（条件不对称性）：}}] 条件概率一般不具有对称性，$P(B \mid A)$与$P(A \mid B)$通常不同。
		\item[\textbf{\textcolor{CTrap}{错因分析（对称错觉）：}}] 学生受“且”的对称性影响，以为条件也可随意交换。
	\end{description}
\end{trapbox}
